\documentclass[12pt,a4paper]{article}
\usepackage[margin=2.5cm]{geometry}
\usepackage[T1]{fontenc}
\usepackage[utf8]{inputenc}
\usepackage[brazil]{babel}
\usepackage{lmodern}
\usepackage{hyperref}
\usepackage{enumitem}
\usepackage{xcolor}

\title{Alento v2.0 \\ Handoff para Nova Branch}
\author{Instituto Alento}
\date{\today}

\begin{document}
\maketitle
\tableofcontents
\newpage

\section{Objetivo deste documento}
Este documento prepara a base para a branch \texttt{v2.0}, onde novas modalidades e novos modulos serao criados.
O foco e permitir evolucao sem quebrar a operacao atual.

\section{Estrategia de branching recomendada}
\begin{itemize}[leftmargin=*]
\item Base atual: \texttt{main} (producao estavel).
\item Nova frente: \texttt{feature/v2.0-core}.
\item Sub-branches por modulo:
\begin{itemize}[leftmargin=*]
\item \texttt{feature/v2.0-modalidades}
\item \texttt{feature/v2.0-modulo-x}
\item \texttt{feature/v2.0-modulo-y}
\end{itemize}
\item Integracao por PR pequeno e frequente para reduzir risco.
\end{itemize}

\section{Contrato minimo para novos modulos}
\begin{enumerate}
\item Definir objetivo funcional do modulo.
\item Definir perfil que pode ver e editar.
\item Definir rotas web e API.
\item Definir schema de dados e validacoes.
\item Definir telas e fluxo de uso.
\item Definir logs e auditoria basica.
\item Definir casos de teste criticos.
\end{enumerate}

\section{Checklist de compatibilidade com v1}
\begin{itemize}[leftmargin=*]
\item Nao quebrar login, sessao e permissoes existentes.
\item Nao remover campos usados por telas ativas.
\item Nao alterar resposta de API sem versao ou adaptacao.
\item Manter rotas antigas funcionando durante transicao.
\item Migracoes de dados devem ser reversiveis quando possivel.
\end{itemize}

\section{Padrao de implementacao para a v2.0}
\subsection{Backend}
\begin{itemize}[leftmargin=*]
\item Controller fino e regra no service.
\item Validacao de entrada antes de consultar o banco.
\item Mensagens de erro claras e padronizadas.
\end{itemize}

\subsection{Frontend}
\begin{itemize}[leftmargin=*]
\item Formularios com validacao em tempo real.
\item Feedback visual simples (erro e sucesso).
\item Acessibilidade minima: foco, contraste e labels.
\end{itemize}

\subsection{Seguranca}
\begin{itemize}[leftmargin=*]
\item Reforcar controle de permissao por rota.
\item Revisar CSRF e rate limit em novos endpoints.
\item Evitar exposicao de dados sensiveis em logs.
\end{itemize}

\section{Plano de migracao sugerido}
\begin{enumerate}
\item Congelar requisitos da v2.0 em backlog fechado.
\item Implementar primeiro a camada de dados e permissoes.
\item Implementar telas em modulos menores.
\item Rodar teste regressivo da v1 a cada merge.
\item Liberar homologacao interna.
\item Liberar producao em lotes.
\end{enumerate}

\section{Criterios para aceite}
\begin{itemize}[leftmargin=*]
\item Build e testes verdes.
\item Fluxos principais da v1 funcionando.
\item Novos fluxos da v2.0 com evidencia de teste.
\item Documentacao atualizada no mesmo PR do modulo.
\end{itemize}

\section{Resumo para equipe}
\textbf{Regra de ouro:} evoluir rapido, mas com isolamento por branch, contrato claro por modulo e integracao curta e frequente.

\end{document}
